% Ad Familiares 3.2 --- headnote
% ad-familiares-03-02, ca. 15 March 51 BC, Rome

\textit{Cicero to Appius Claudius Pulcher, proconsul of Cilicia,
written at Rome around the middle of March 51 BC (Perseus:
\textit{Romae circ. m. Mart. a. 703}), as Cicero prepared to leave
the city for the province Appius had been holding for the past
two years. The lex Pompeia of 52 BC had imposed a five-year
interval between magistracy and promagistracy, and the unwilling
consular --- who had spent his career steering clear of
provincial commands --- was caught in the new dispensation and
sent out to Cilicia in Appius's place. The handover was, by all
reports of the months that followed, a disaster: Appius was
draining the province and intended to draw out his stay; Cicero,
once on the ground, would find decaying garrisons and an
exhausted treasury. None of this can be said. The Pulchri are
the most dangerous family in Rome, and Appius is brother to
Publius Clodius; nothing the incoming proconsul writes can be
allowed to look like complaint.}

\textit{What he writes instead is the polished superficial cover
that the situation requires. He casts himself as a reluctant
successor, comforted only by the thought that no one is more
Appius's friend than he and that no one therefore would be more
likely to receive the province in good order from him. The
substantive request --- the heart of the letter --- is buried in
two sentences of formula: that Appius, by every means in his
power (and many will be in his power), look ahead and consult
for Cicero's interests; that, so far as he can, he hand over the
province in the freest possible condition. The opener
\textit{M. Cicero procos. s. d. Appio Pulchro imp.} announces
the formal symmetry of the transaction: two commanders, two
titles, equal and courteous. The whole short letter is the
careful first move in a correspondence whose two participants
both know that the substance is going to be harder than the
form.}
